\section{Exact positive cover counts for a two-prefix profile}

\subsection{O-026: a fixed-base optimization within the existing criteria}
Let $L\geq M\geq b\geq1$, let $r\geq4$, and put
\[
 q=r-3,\qquad u=L-b,\qquad v=M-b,\qquad
 c=(v+1)(u+1)-\frac{v(v+1)}2.
\]
Consider the support
\[
 \Gamma=\left\{\gamma:
       |\gamma|\leq L,\quad
       \sum_{i=3}^r\gamma_i\leq M,\quad
       \sum_{i=4}^r\gamma_i\leq b\right\}.
                                                               \tag{O-026.1}
\]
Assume
\[
                             q>\max\{c,u+v+2\}.
                                                               \tag{O-026.2}
\]
Among nonempty finite monomial order ideals, the following
two exact minima hold:
\[
 \begin{split}
 \min_{\Gamma\subseteq A+D_b}|A|&=c,\\
 \min_{\Gamma\subseteq A+B+D_b}(|A|+|B|)&=u+v+2.
 \end{split}                                                  \tag{O-026.3}
\]
The unique one-sided optimum is
\[
                  A_0=\{x_2^a x_3^t:a+t\leq u,\ t\leq v\}.
                                                               \tag{O-026.4}
\]
Two axis ideals of degrees $u$ and $v$ attain the
two-sided minimum. These are exact positive counts for
the fixed-base support optimization in M-018 and M-019.
They are consequences and optimality statements within
those already proved criteria, not new existential
multiplicity criteria.

\begin{proof}[The one-sided optimum]
O-024 gives $\Gamma\subseteq A_0+D_b$, and its
subpartition count gives $|A_0|=c$.
Take any proposed order ideal $A$ with $|A|\leq c$.
Since there are $q>c$ suffix variables, some unit
$x_j$, $j\geq4$, is absent from $A$.
Downward closure implies that every member of $A$
has zero $j$-th exponent.

For every prefix monomial $\nu\in A_0$,
$x_j^b\nu$ belongs to $\Gamma$:
its total degree is at most $b+u=L$, its suffix from
coordinate three is at most $b+v=M$, and its final
suffix degree is $b$. In a decomposition of its exponent
as $\alpha+\eta$, with $\alpha\in A$ and $|\eta|\leq b$,
the $j$-th exponent forces $\eta=be_j$.
Thus $\alpha=\nu$, proving $A_0\subseteq A$.
Therefore $|A|\geq c$; equality forces $A=A_0$.
Any ideal with more than $c$ members cannot improve
the known optimum. This proves both the minimum and
its uniqueness.
\end{proof}

\begin{proof}[The two-sided optimum]
The two axis ideals
$A=\{1,x_2,\ldots,x_2^u\}$ and
$B=\{1,x_3,\ldots,x_3^v\}$ contain a sum covering
$A_0$, so they attain the upper bound $u+v+2$.
Now take any proposed pair with
$|A|+|B|\leq u+v+2$.
By (O-026.2) there is a suffix unit $x_j$ in neither
ideal, and no member of either ideal uses that variable.

Apply the cover to $x_j^b x_2^u$.
Its $D_b$ part must be precisely $x_j^b$, so its
remaining decomposition uses pure $x_2$ powers in
$A$ and $B$. If $a_2,b_2$ are their largest $x_2$
axis degrees, this forces $a_2+b_2\geq u$.
Likewise $x_j^b x_3^v$ gives $a_3+b_3\geq v$.
Every ideal contains its two full axis chains, sharing
only the constant monomial. Consequently
\[
 |A|+|B|\geq(a_2+a_3+1)+(b_2+b_3+1)
                              \geq u+v+2.
\]
Larger pairs cannot improve the already attained upper
bound, proving the second minimum.
\end{proof}

\paragraph{The exact counts in the positive example.}
For O-025, take
$L=1000006$, $M=1000000$, and $b=26$.
Then
\[
 \begin{split}
 u&=999980,\qquad v=999974,\\
 c&=499981500150,\qquad u+v+2=1999956.
 \end{split}                                                  \tag{O-026.5}
\]
The suffix count $r-3=2^{400}-3$ satisfies (O-026.2).
At the common fixed core size
$T=m\binom{r+11}{12}$, the optimal one-sided
matrix size is therefore $499981500151T$, whereas
the optimal two-sided size at this same base order is
$1999956T$.
O-025 has already proved that the latter fits inside
$n$ and the former does not, with actual permanent
multiplicity at least four.

For this particular example the one-sided count can
also be optimized over every even base order of M-018.
The optimum still occurs at $b=26$ with the unique
ideal (O-026.4). Indeed its size is less than
$2^{39}n$, since $T<n$ and
$499981500151<2^{39}$.
For $b\geq28$, (O-025.6) gives a lower ratio to $n$
of at least $2r/13!>2^{39}$.
For $b\leq24$, (O-025.9) gives a lower ratio greater
than $r/2^{139}=2^{261}>2^{39}$.
These cases exhaust the other orders. Enlarging the
ambient length cannot lower the core size or help
decompose monomials supported on the original variables.
This establishes the stated optimum for the M-018
architecture, without an optimality assertion for
arbitrary determinant maps or for all M-019 base orders.

\paragraph{Geometric dictionary.}
The counts in (O-026.3) are ranks of explicitly indexed
free coefficient modules over $\mathbb Z$, with their
torus-comodule structures from SGA~3 I,
Proposition~4.7.3, pp.~35--36.
Nonnegative character exponents make the missing-variable
argument exact: a monomial containing that character
cannot be cancelled inside a support decomposition.
The already proved morphisms M-018 and M-019 convert
the counts into universally injective coordinate
pullbacks. Their characteristic-zero Schur-copy step
uses SGA~6 VI, Theorem~3.3, pp.~392--393, and the
Cauchy/Schur formalism of Macdonald I(4.3), p.~63,
Appendix~A(5.3)--(5.6), pp.~155--156.
The optimization concerns these explicit coefficient
schemes and maps; neither a virtual lambda-ring identity
nor purity alone provides the comparison.
